\section{Introdu\c{c}\~ao}

\subsection{Contexto: a linhagem do c\'odigo}

O HiGFlow resolve as equa\c{c}\~oes de Navier--Stokes incompress\'iveis por
diferen\c{c}as finitas sobre malhas em \'arvore n\~ao graduadas, com
interpola\c{c}\~ao sem malha nas interfaces de refinamento~\cite{sousa2019}. Ele
herda uma tradi\c{c}\~ao longa de simula\c{c}\~ao de escoamentos com superf\'icie
livre no grupo: o sistema Freeflow~\cite{castelo2000}, constru\'ido sobre o
m\'etodo GENSMAC~\cite{tome1994} e sua extens\~ao tridimensional, e a
generaliza\c{c}\~ao para escoamentos multif\'asicos por
\emph{front-tracking}~\cite{sousa2004}.

Essa \'ultima refer\^encia \'e a mais pr\'oxima deste trabalho, e vale dizer
explicitamente em que: \cite{sousa2004} implementa front-tracking sobre o
GENSMAC-3D, com a frente representada por uma malha de superf\'icie
lagrangeana acoplada a uma malha euleriana de diferen\c{c}as finitas. O que este
relat\'orio documenta \'e a constru\c{c}\~ao do an\'alogo dentro do HiGFlow,
come\c{c}ando em 2D, sobre a maquinaria de fronteira imersa que j\'a existia no
sistema.

\subsection{Front-tracking \'e alternativa ao VOF, n\~ao complemento}

Os dois m\'etodos resolvem o mesmo problema f\'isico por representa\c{c}\~oes
opostas da interface:

\begin{itemize}
  \item \textbf{VOF} (\emph{volume of fluid}) \emph{captura} a interface por um
    campo euleriano de fra\c{c}\~ao volum\'etrica advectado com o escoamento. A
    interface n\~ao existe como objeto; \'e reconstru\'ida da fra\c{c}\~ao quando
    preciso. A tens\~ao superficial entra tipicamente pela formula\c{c}\~ao de
    for\c{c}a de superf\'icie cont\'inua (CSF)~\cite{brackbill1992}.
  \item \textbf{Front-tracking} \emph{rastreia} a interface por marcadores
    lagrangeanos conectados que se movem com o fluido~\cite{tryggvason2001}. A
    interface \'e um objeto geom\'etrico expl\'icito, com conectividade, e a
    curvatura sai dessa geometria.
\end{itemize}

A decis\~ao de projeto, tomada no in\'icio e mantida, \'e que os dois ficam
\textbf{separados no c\'odigo}: n\~ao se mistura a reconstru\c{c}\~ao geom\'etrica
de um com a advec\c{c}\~ao de fra\c{c}\~ao do outro. \'E o que torna poss\'ivel a
compara\c{c}\~ao da Se\c{c}\~ao~\ref{sec:comparacao}: h\'a dois m\'etodos
independentes no mesmo sistema, sobre o mesmo solver.

\subsection{Escopo: o que est\'a feito}

O projeto est\'a organizado em fases, cada uma com um or\'aculo exato antes de a
pr\'oxima introduzir f\'isica nova. O estado no momento deste relat\'orio:

\begin{center}
\begin{adjustbox}{max width=\linewidth}
\begin{tabular}{llp{7.2cm}}
\toprule
Fase & Estado & Or\'aculo \\
\midrule
B1 & feita & advec\c{c}\~ao cinem\'atica reversivel de Rider--Kothe: a frente
             volta \`a forma inicial em $t=T$ \\
B2 & feita & gota est\'atica: o salto de press\~ao bate com a lei de Laplace
             $\Delta p = \sigma/R$, e as correntes par\'asitas ficam abaixo de um
             teto medido \\
B3 & aberta & oscila\c{c}\~ao de gota na frequ\^encia de Lamb \\
B4 & aberta & bolha subindo, \emph{benchmark} publicado \\
B5 & aberta & 3D: superf\'icie triangulada \\
\bottomrule
\end{tabular}
\end{adjustbox}
\end{center}

A ordem \'e deliberada: B1 e B2 n\~ao t\^em din\^amica acoplada completa, e \'e
neles que os defeitos de m\'aquina --- advec\c{c}\~ao, cirurgia, curvatura,
espalhamento --- aparecem baratos. Tudo que se segue foi medido; onde um teste
n\~ao exercita algo, isso est\'a dito.
